% Sorted in alphabetical order. 

\begin{description}[style=nextline]

\item[Anti-Rigid Type] A type that classifies each of its instances contingently: every instance of the type can cease to instantiate it without ceasing to exist.

\item[Conceptualization] A system of concepts and conceptual relations in the mind of an agent, or shared by a community of agents, through which a domain is perceived, organized, and understood. A conceptualization is not itself a model artifact; rather, it is what a conceptual model aims to explicitly describe or approximate.

\item[Foundational Ontology] A domain-independent ontology that defines highly general categories and relations, used as a semantic foundation for more specific domain ontologies and modeling languages.

\item[Identity Principle] The principle by which an individual is distinguished from all others and can be recognized as the same individual across change.

%\item[Identity Principle] A principle used to determine whether two occurrences, references, or temporal manifestations correspond to the same individual.

\item[Identity Provider] A sortal type that supplies an identity principle to its instances.

\item[Non-Sortal Type] A type that does not provide or carry a uniform identity principle for all of its instances, classifying instances that follow different identity principles.

\item[Ontology] A reference model of consensus within a community, formally defining the types of entities that may exist in a domain, the properties they may bear, the relations in which they may participate, and the constraints governing them.

\item[Ontology-driven Conceptual Modeling] A conceptual modeling approach in which modeling languages, modeling constructs, and modeling decisions are grounded in well-founded ontological theories. Its purpose is to produce conceptual models that make explicit the ontological commitments of a domain and support clear communication, conceptual analysis, and semantic interoperability.

\item[Rigidity] A modal meta-property of a type indicating whether instances of the type instantiate it necessarily (i.e., in all possible circumstances) or contingently (i.e., in some circumstances but not in all of them).

\item[Rigid Type] A type that classifies each of its instances essentially: every instance of the type must instantiate it in every possible circumstance in which that instance exists.

\item[Semi-Rigid Type] A type that essentially classifies some of its instances and contingently classifies others. 

\item[Sortality] A meta-property of a type that concerns whether the type classifies instances that share a uniform principle of identity.

\item[Sortal Type] A type whose instances follow a uniform identity principle.




\end{description}